\documentclass[11pt]{amsart}

\usepackage{amsmath,amssymb,amsthm,mathtools}
\usepackage[a4paper,margin=26mm]{geometry}

\newtheorem{theorem}{Theorem}[section]
\newtheorem{lemma}[theorem]{Lemma}
\newtheorem{corollary}[theorem]{Corollary}
\theoremstyle{remark}
\newtheorem{remark}[theorem]{Remark}

\newcommand{\Z}{\mathbb Z}
\newcommand{\Q}{\mathbb Q}
\newcommand{\aug}{\varpi}
\newcommand{\gm}{\gamma}

\title{A short two-factor proof for metabelian Lie rings}
\author{}
\date{}

\begin{document}

\begin{abstract}
We prove the stronger nilpotent form of the odd-dimensional inclusion:
if a metabelian Lie ring $K$ has nilpotency class at most $c\geq2$, then
$\delta_{2c-1}(K)=0$.  The proof uses one PBW functional supported on
monomials with at most two Lie factors.  A Smith ordering makes every
higher relation row invisible to this functional; its only possible
correction is a central quadratic commutator.  Thus no reduction to a
finite ring or a cyclic last layer, and no packet calculation, is needed.
\end{abstract}

\maketitle

For a Lie ring $A$, write $U(A)$ for its universal enveloping algebra,
$\aug(A)$ for its augmentation ideal, and
\[
  \delta_s(A)=A\cap\aug(A)^s,
  \qquad \gm_{s+1}(A)=[\gm_s(A),A].
\]
All Lie rings are over $\Z$.

\section{The ordered two-factor functional}

We first isolate the only presentation calculation that is needed.

\begin{lemma}[central Smith corrections]\label{lem:central}
Let $F$ be relatively free of finite rank in the variety of metabelian
Lie rings of class at most $c$, and write
\[
                         F=X\oplus M,\qquad M=F'.
\]
Thus $M$ is abelian.  Let $R\mathrel{\unlhd}F$, put $K=F/R$, and let
$E\subseteq X$ be the image of $R$ under $F\to X=F/M$.  There is a
basis $x_1,\ldots,x_m$ of $X$, an integer $r\leq m$, and integers
\[
                    0<d_1\mid d_2\mid\cdots\mid d_r
\]
such that $E$ has basis $d_1x_1,\ldots,d_rx_r$.  Choose
\[
                  \rho_i=d_ix_i+b_i\in R,
                  \qquad b_i\in M\quad(1\leq i\leq r).
\]
For $1\leq j<i\leq r$, put
\[
                         q_{ij}=d_i[x_i,x_j].
\]
Then the images of all $q_{ij}$ are central in $K$.
\end{lemma}

\begin{proof}
We calculate in $K$ and omit bars.  The relations give
\[
        d_ix_i=-b_i,
        \qquad d_ix_j=-\frac{d_i}{d_j}b_j\quad(j<i).
\]
For every generator $x_k$, commutativity of the adjoint $X$-actions
on the abelian ideal $K'$ gives
\begin{align*}
 [q_{ij},x_k]
 &=-[b_i,x_j,x_k]
  =-[b_i,x_k,x_j]                                      \\
 &=d_i[x_i,x_k,x_j]
  =[[x_i,x_k],d_ix_j]=0.
\end{align*}
The last bracket is zero because both its entries belong to $K'$.
Also $[q_{ij},K']=0$.  Hence $q_{ij}$ is central.
\end{proof}

Let $C$ be any central subgroup of $K$, and let $Q$ be the subgroup
generated by the images of the $q_{ij}$.  By
Lemma~\ref{lem:central},
\[
                              Z=C+Q
\]
is central.  Fix an additive character $\ell:K\to\Q/\Z$ and put
$W=K/Z$.  The rule
\[
 \Omega(\bar x,\bar y)=\ell([x,y])                    \tag{1}
\]
is a well-defined alternating biadditive form on $W$.

There is a biadditive $H:W\times W\to\Q/\Z$ such that
\[
 H(u,v)-H(v,u)=\Omega(u,v).                            \tag{2}
\]
Indeed, decompose the finitely generated abelian group $W$ into cyclic
summands, order their generators, put the matrix of $H$ strictly above
the diagonal equal to that of $\Omega$, and put all other entries equal
to zero.  The orders of the cyclic generators also annihilate the
corresponding entries of $\Omega$, so this prescription is well defined.

Choose a homogeneous basis of $M$ and place all its elements before
\[
                         x_1<x_2<\cdots<x_m              \tag{3}
\]
in a PBW order for $F$.  Define an additive functional
$\Lambda:U(F)\to\Q/\Z$ on the ordered PBW monomials by
\[
 \begin{array}{c|c}
 \text{number of Lie factors}&\Lambda\\ \hline
 1&\ell(\text{its image in }K),\\
 2&H(\text{image of the first factor},
       \text{image of the second factor}),\\
 0\text{ or at least }3&0.
 \end{array}                                             \tag{4}
\]
In the two-factor line the images are taken in $W$.

The relation $uv-vu=[u,v]$ and (2) immediately imply
\begin{equation}
             \Lambda(uv)=H(\bar u,\bar v)
             \qquad(u,v\in F).                          \tag{5}
\end{equation}
Here the product on the left is collected in the fixed PBW order,
whereas the arguments on the right retain their displayed order.
Since $R$ is a Lie ideal, $x\rho=\rho x+[x,\rho]$ shows that the
two-sided ideal of $U(F)$ generated by $R$ is exactly $RU(F)$.

\begin{lemma}[relation rows are invisible]\label{lem:rows}
With the preceding notation,
\[
                              \Lambda(RU(F))=0.
\]
\end{lemma}

\begin{proof}
Put $N=R\cap M$.  Additively,
\[
                     R=N+\sum_{i=1}^r\Z\rho_i.           \tag{6}
\]
It is therefore enough to multiply an element of $N$, or one of the
$\rho_i$, by an ordered PBW monomial $u$.

Write the ordered multiplier as
\[
 u=m_1\cdots m_sx_{j_1}\cdots x_{j_t},
 \qquad j_1\leq\cdots\leq j_t,
 \qquad k=s+t,
\]
where the $m_a$ are basis elements of $M$.  If $n\in N$, inserting $n$ into the
initial $M$-block of $u$ creates no commutator, since $M$ is abelian.
For $k=0$ its value is $\ell(n)=0$; for $k=1$ equation (5) gives
$H(0,\bar u)=0$; and for $k\geq2$ the collected products have at least
three factors.  Thus $\Lambda(nu)=0$.  More generally, if an
$N$-factor occurs at an intermediate position during collection, moving
it into the $M$-block can only create further $N$-factors, because $N$
is an ideal.  Hence every terminal term with at most two factors still
has an $N$-factor and has value zero by the same argument.

Consider $\rho_i u=(d_ix_i+b_i)u$.  For $k=0$ the assertion is
$\ell(\rho_i)=0$, and for $k=1$ it follows from (5).  Suppose
$k\geq2$.  The term $b_i u$ is collected merely by inserting $b_i$
into the initial abelian $M$-block, so it retains $k+1$ factors and is
invisible to $\Lambda$.

It remains to insert $d_ix_i$ into $u$.  When it crosses a factor
$m\in M$, its correction is
\[
 d_i[x_i,m]=[\rho_i,m]\in N,                             \tag{7}
\]
because $[b_i,m]=0$.  After the $M$-block has been crossed, the only
ordinary factors which $x_i$ must cross are $x_j$ with $j<i$.  The
correction is then precisely $q_{ij}$.  If this derived correction has
to move left past an earlier $x_j$, every further correction belongs
to $N$, because $[q_{ij},F]\subseteq R$ by
Lemma~\ref{lem:central}.  A branch in which $q_{ij}$ is never bracketed
retains exactly $k$ factors; once it is bracketed, it becomes an
$N$-branch.  Consequently exact adjacent PBW collection
has only the following three kinds of terminal terms:
\[
 \begin{array}{ll}
 \text{(a)}&\text{the principal insertion, with }k+1\text{ factors};\\
 \text{(b)}&\text{a term with a factor in }N;\\
 \text{(c)}&\text{a term with a factor }q_{ij},
              \text{ with exactly }k\text{ factors}.
 \end{array}                                             \tag{8}
\]
Terms of type (a) vanish by (4), and terms of type (b) vanish by the
first paragraph.  A term of type (c) has at least two factors.  It is
zero under $\Lambda$ if it has at least three; if it has two, (5)
applies and one argument is zero in $W$, since $q_{ij}\in Z$.
This proves $\Lambda(\rho_i u)=0$ and hence the lemma.
Notice that the split $\rho_i=d_ix_i+b_i$ was used only for
$k\geq2$, where the two factor-preserving parts are separately invisible.
For $k=0,1$, equation (5) evaluates the full relation.  Thus no tail
$b_i$ has been treated as a relation.
\end{proof}

\section{The weight cutoff}

\begin{theorem}\label{thm:nilpotent}
Let $c\geq2$.  If $K$ is a metabelian Lie ring of nilpotency class at
most $c$, then
\[
                              \delta_{2c-1}(K)=0.
\]
\end{theorem}

\begin{proof}
It is enough to treat finitely generated $K$.  Indeed, from a free
presentation $K=\widehat F/\widehat R$ and a nonzero
$a\in\delta_{2c-1}(K)$ choose a lift and an equality
\[
 \widetilde a=\sum_{\nu=1}^t\rho_\nu u_\nu+v,
 \qquad \rho_\nu\in\widehat R,
 \quad v\in\aug(\widehat F)^{2c-1}.                       \tag{9}
\]
Only finitely many free generators occur here.  The relatively free
subring $F_0$ on those generators is a retract of $\widehat F$, so
(9) holds in $U(F_0)$.  If $R_0$ is the ideal of $F_0$ generated by
the $\rho_\nu$, then $K_0=F_0/R_0$ maps to $K$; the image $a_0$ of
$\widetilde a$ belongs to $\delta_{2c-1}(K_0)$ and maps to $a$.
Thus $a_0\ne0$, reducing the question to the finitely generated ring
$K_0$.

Present $K=F/R$, where $F$ is the relatively free metabelian Lie ring
of class at most $c$ on finitely many generators.  It has the usual
positive grading
\[
             F=F_1\oplus\cdots\oplus F_c,
             \qquad F_1=X,
\]
whose homogeneous groups are free abelian.  Use homogeneous bases in
the PBW order (3).  If a PBW monomial is assigned the sum of the
weights of its Lie factors, then
\begin{equation}
          \aug(F)^s=
          \langle\text{PBW monomials of total weight at least }s\rangle.
                                                               \tag{10}
\end{equation}
Indeed, $F_t\subseteq\aug(F)^t$, while a product of $s$ positive-weight
elements has weight at least $s$.

Suppose that $0\ne a\in\delta_{2c-1}(K)$.  Choose an additive character
$\ell:K\to\Q/\Z$ with $\ell(a)\ne0$; such characters separate points
of every finitely generated abelian group, as follows immediately from
its decomposition into cyclic summands.  In the construction of
Section~1 take
\[
                              C=\gm_c(K).
\]
Choose a lift $\widetilde a\in F$.  Membership in the dimension
subring gives
\begin{equation}
             \widetilde a\in RU(F)+\aug(F)^{2c-1}.       \tag{11}
\end{equation}

Lemma~\ref{lem:rows} says that $\Lambda$ vanishes on the first summand
of (11).  It also vanishes on the second.  To see this, use (10).  A
one-factor PBW monomial has weight at most $c<2c-1$.  If a two-factor
monomial has weight at least $2c-1$, one of its factors has weight $c$,
because two weights at most $c-1$ have sum at most $2c-2$.  That factor
maps into $\gm_c(K)=C$ and hence to zero in $W$.  Monomials with at
least three factors vanish by the definition of $\Lambda$.

Thus (11) gives
\[
                       0=\Lambda(\widetilde a)=\ell(a),
\]
contrary to the choice of $\ell$.  Hence $\delta_{2c-1}(K)=0$.
\end{proof}

\begin{corollary}\label{cor:main}
For every metabelian Lie ring $L$ and every $n\geq1$,
\[
                         \delta_{2n+1}(L)
                         \subseteq\gm_{n+2}(L).
\]
\end{corollary}

\begin{proof}
Apply Theorem~\ref{thm:nilpotent}, with $c=n+1$, to
$K=L/\gm_{n+2}(L)$.  The quotient map sends
$\delta_{2n+1}(L)$ into $\delta_{2n+1}(K)=0$.
\end{proof}

For $n=2$ the same conclusion holds for an arbitrary Lie ring $L$:
$L/\gm_4(L)$ has class at most three and is therefore metabelian.

\begin{remark}
The number $2c-1$ enters only once: it is the first weight at which a
two-factor PBW monomial is forced to contain a factor from the last
central layer.  The metabelian hypothesis enters only in the ordered
row lemma: it makes the derived block abelian and makes every Smith
swap correction $q_{ij}$ central.  This is why the calculation stops
at two factors.
\end{remark}

\end{document}
